% !TeX root = ../初等数学笔记.tex
% 由主文档通过 \input 载入；请编译主文档。
\setcounter{chapter}{1} % 保留原书章号。
\chapter{数系、运算与近似}
\label{chap:em-numbers}
数系的扩展来自运算和测量的需要：减法引入负数，除法引入分数，几何长度进一步要求无理数。本章把实数的四则运算、顺序和完备性作为基础结构，重点研究正确运算与可靠近似；实数的构造属于数学分析的基础问题。

\section{整数、分数与有理数}
\emph{有理数}（rational number）是可写成 $p/q$ 的数，其中 $p,q\in\mathbb Z$、$q\ne0$。分数的不同表示满足
\[
 \frac pq=\frac rs\Longleftrightarrow ps=rq,\qquad
 \frac pq+\frac rs=\frac{ps+rq}{qs},\qquad
 \frac pq\frac rs=\frac{pr}{qs}.
\]
约分改变表示而不改变数值，但只能约去共同的非零因子。除法是乘以倒数，除数为零时没有定义，因为 $0\cdot x=b$ 在 $b\ne0$ 时无解，在 $b=0$ 时也不能唯一确定 $x$。

比例比较必须统一单位。若某比例由 $20\%$ 增至 $25\%$，其增加量为 $5$ 个百分点，而相对增加为 $(25-20)/20=25\%$。若商品先涨价 $10\%$ 再降价 $10\%$，最终价格乘以 $1.1\times0.9=0.99$，比原价低 $1\%$，两次百分比作用的基数不同。

\section{无理数与实数轴}
不属于 $\mathbb Q$ 的实数称为\emph{无理数}（irrational number）。以 $\sqrt2$ 为例，若 $\sqrt2=p/q$ 且 $p,q$ 已约至互素，则 $p^2=2q^2$。平方为偶数蕴含本身为偶数，故 $p=2k$；代入得 $q^2=2k^2$，又推出 $q$ 为偶数，与互素矛盾。最简分数的存在可由正分母取最小值得到，不需要预先使用素因数分解。

有理数的小数表示有限或最终循环：长除法中的余数只有有限种，余数为零则终止，否则必有余数重复。反之，最终循环小数可通过乘以适当的 $10$ 的幂再相减转化为分数，例如 $x=0.\overline{27}$ 满足 $100x-x=27$，故 $x=3/11$。小数表示本身未必唯一，如 $0.999\ldots=1$；无限小数的值需要通过极限或实数完备性解释。

\emph{实数}（real number）与数轴上的点相对应，$\mathbb R=\mathbb Q\cup(\mathbb R\setminus\mathbb Q)$。本书采用实数的完备性：每个非空且有上界的实数集合有最小上界。它保证数轴没有有理数系中那样的空缺；该性质并非仅由四则运算律推出。

\section{顺序、绝对值与距离}
实数的顺序与加法相容；乘以正数保持大小关系，乘以负数反向。定义
\[
 |x|=\begin{cases}x,&x\ge0,\\-x,&x<0.\end{cases}
\]
于是 $|x|\ge0$、$|xy|=|x||y|$，且 $-|x|\le x\le|x|$。两数 $x,y$ 在数轴上的距离为 $|x-y|$。对 $r>0$，$|x-a|<r$ 等价于 $a-r<x<a+r$，$|x-a|\ge r$ 等价于 $x\le a-r$ 或 $x\ge a+r$。

\begin{theorem}[三角不等式与反三角不等式]
对任意实数 $x,y$，
\[
 |x+y|\le |x|+|y|,\qquad
 \bigl||x|-|y|\bigr|\le |x-y|.
\]
第一式等号成立当且仅当 $xy\ge0$。
\end{theorem}
\begin{proof}
将 $-|x|\le x\le|x|$ 与对应的 $y$ 不等式相加，得到第一式。两侧非负，其平方之差为 $2(|xy|-xy)$，故等号条件为 $xy\ge0$。再由 $x=(x-y)+y$ 得 $|x|-|y|\le|x-y|$，交换 $x,y$ 即得第二式。
\end{proof}

\section{整数幂、根式与运算范围}
正整数幂由重复相乘定义；当 $a\ne0$ 时，规定 $a^0=1$、$a^{-n}=1/a^n$。指数律 $a^ma^n=a^{m+n}$ 与 $(a^m)^n=a^{mn}$ 在相应运算有定义时成立。本书不把 $0^0$ 作为通常实数幂定义。

当 $a\ge0$ 时，$\sqrt a$ 是平方等于 $a$ 的唯一非负实数，因此 $\sqrt{x^2}=|x|$。偶次实根要求被开方数非负并取非负值，奇次实根对所有实数存在且唯一。对 $a,b\ge0$ 有 $\sqrt{ab}=\sqrt a\sqrt b$，因为右侧非负且平方为 $ab$；若不满足范围条件，不能照搬公式。

有理化利用平方差，例如 $a,b\ge0$ 且 $a\ne b$ 时，
\[
 \frac1{\sqrt a-\sqrt b}=\frac{\sqrt a+\sqrt b}{a-b}.
\]
它既可用于化简，也能避免两个接近的近似数相减所造成的相对误差放大。

\section{近似数、有效数字与误差}
设准确值为 $x$，近似值为 $\widetilde x$。绝对误差为 $|\widetilde x-x|$，在 $x\ne0$ 时相对误差为 $|\widetilde x-x|/|x|$。保留小数点后 $m$ 位并四舍五入时，绝对误差至多为 $\frac12\,10^{-m}$；有效数字则从首个非零数字起计，反映相对精度的数量级。

若 $|x-\widetilde x|\le\varepsilon$、$|y-\widetilde y|\le\delta$，则三角不等式给出
\[
 |(x+y)-(\widetilde x+\widetilde y)|\le\varepsilon+\delta,
 \quad
 |xy-\widetilde x\widetilde y|
 \le|\widetilde x|\delta+|\widetilde y|\varepsilon+\varepsilon\delta.
\]
第二式由 $x=\widetilde x+u,y=\widetilde y+v$ 展开得到。商的误差控制还要求分母远离零；仅知道分母误差很小并不足够。计算通常保留若干保护位，在最终结果处统一舍入。

\section{例题、习题与学习检查}
\begin{example}
不使用小数近似，化简 $\sqrt{(x-2)^2}$；再估计 $\sqrt2$ 到三位小数。
\end{example}
\begin{solve}
第一式等于 $|x-2|$，即 $x\ge2$ 时为 $x-2$，$x<2$ 时为 $2-x$。又 $1.414^2=1.999396<2$，$1.415^2=2.002225>2$，故 $1.414<\sqrt2<1.415$。进一步有 $1.4145^2=2.00081025>2$，所以四舍五入到三位小数为 $1.414$，仅有前两个界还不能确定舍入方向。
\end{solve}
\begin{enumerate}
\item 化 $0.1\overline6$ 为分数。答：$1/6$。
\item 举例说明两无理数之和、之积都可能为有理数。可取 $\sqrt2,-\sqrt2$，其和为 $0$、积为 $-2$。
\item 长方形两边近似为 $3,4$，各有绝对误差不超过 $0.01$，求面积误差界。答：相对于近似面积 $12$，误差不超过 $0.0701$。
\end{enumerate}

